\section{Exact finite-model results}
\label{sec:closed-loops}

The abstract influence-algebra construction identifies the kind of
information that may survive coarse graining, but it does not by itself
provide a microscopic dilation, a mixing rate, a phase transition, or a
geometric estimator.  This section records independent models that supply
these missing pieces.  They are presented together because they
constrain the architecture of RQCP-QG; they are not combined into one
generator.  In particular, a solid implication below is always local to the
stated model.

\subsection{Local unitaries to quantum-response kinetics}
\label{sec:loop-unitary-kinetics}

\paragraph{Assumption removed.}
The original process description began with a UCP coarse-graining channel.
That is too high a starting point for a microscopic theory: complete
positivity may be postulated while the Markov property and the kinetic rates
remain unexplained.  A collision circuit supplies an exact dilation of
the channel used in the Boolean sector model
\cite{XuResponseKinetics2026}.  The continuum generator belongs to the
standard Gorini--Kossakowski--Sudarshan--Lindblad class
\cite{Gorini1976,Lindblad1976}; the point here is that its rates and its
finite-step error follow from an explicit one-pass quantum network, rather
than being postulated at the semigroup level \cite{Chiribella2009}.

For an involution \(U=U^\dagger=U^{-1}\), let a system interact once with a
fresh qubit bath cell through
\begin{equation}
 V_\theta=\exp(-i\theta\,U\otimes X_b),
 \qquad
 \rho_b=\tfrac12\mathbf1_b .
 \label{eq:loop-collision-unitary}
\end{equation}
Tracing out the cell gives
\begin{equation}
 \Phi_{\theta,U}(X)
 =\cos^2\theta\,X+\sin^2\theta\,UXU
 =\exp\!\left[h r(\operatorname{Ad}_U-\operatorname{id})\right]X
 \label{eq:loop-one-collision}
\end{equation}
when \(\sin^2\theta=(1-e^{-2rh})/2\).  The equality is finite-step and
contains no weak-coupling expansion.  Taking \(U=Z_i\) produces local
dephasing; taking \(U=S_{ij}\), the SWAP on an edge, produces symmetric
exchange.  An edge coloring of a bounded-degree graph arranges the exchange
collisions into a constant-depth symmetric round.  Every collision uses a
new cell, so the global reduced circuit is exactly CPTP at every step.

Let \(G=(V,E)\) be connected and \(d\)-regular, and write
\begin{align}
 \mathcal D&=\kappa\sum_i(\operatorname{Ad}_{Z_i}-\operatorname{id}),
 \nonumber\\
 \mathcal A&=\frac{\gamma}{d}\sum_{\{i,j\}\in E}
   (\operatorname{Ad}_{S_{ij}}-\operatorname{id}),
 \qquad
 \mathcal L_{\rm kin}=\mathcal D+\mathcal A .
 \label{eq:loop-kinetic-generator}
\end{align}
For a Pauli string \(P_{\boldsymbol\sigma}\), let \(k\) be its number of
nonidentity labels and \(q\) its number of \(X\) or \(Y\) labels.  The
Heisenberg generator closes exactly:
\begin{equation}
 \mathcal L_{\rm kin}P_{\boldsymbol\sigma}
 =-2\kappa q\,P_{\boldsymbol\sigma}
 +\frac{\gamma}{d}\sum_{\{i,j\}\in E}
  \left(P_{s_{ij}\boldsymbol\sigma}-P_{\boldsymbol\sigma}\right).
 \label{eq:loop-pauli-hierarchy}
\end{equation}
The second term is a colored interchange process and the first kills every
transverse label.  Thus the hierarchy is not a low-moment closure: all Pauli
orders are invariant sectors and together span the full operator algebra.

\begin{theorem}[Finite-step charge algebra and fixed-order kinetic limit]
\label{thm:loop-unitary-kinetics}
Let \(\Phi_h\) be the symmetric collision round constructed from
Eqs.~\eqref{eq:loop-collision-unitary}--\eqref{eq:loop-kinetic-generator}.
For every \(h,\kappa,\gamma>0\),
\begin{equation}
 \operatorname{Fix}(\Phi_h)
 =\operatorname{span}\{P_0,\ldots,P_N\},
 \qquad
 P_m=\sum_{|x|=m}|x\rangle\langle x|.
 \label{eq:loop-finite-fixed}
\end{equation}
On an order-\(k\) Pauli orbit and for coefficient
\(\ell^p\) norm, \(p\in\{1,\infty\}\), \(t=nh\),
\begin{align}
 \|\Phi_h^n-e^{t\mathcal L_{\rm kin}}\|_{(k),p}
 &\le \varepsilon_k(t,h),\nonumber\\
 \varepsilon_k(t,h)
 &=\frac{t h^2}{3}B_k^3e^{hB_k},
 \qquad
 B_k\le\frac{2\chi k\gamma}{d},
 \label{eq:loop-kinetic-bound}
\end{align}
where \(\chi\) is the edge-chromatic number of the chosen schedule.  At
fixed \(k\) the bound is independent of \(N\).
\end{theorem}

\begin{proof}[Proof outline]
Each random-unitary factor is a positive Hilbert--Schmidt contraction.
Equality in the product contraction forces simultaneous commutation with
all \(Z_i\) and all edge SWAPs.  The first condition makes an observable
diagonal, while connected edge transpositions make its diagonal value
depend only on Hamming weight, proving
Eq.~\eqref{eq:loop-finite-fixed}.  On an order-\(k\) orbit, a matching has at
most \(k\) active edges.  The symmetric product and
\(e^{h\mathcal A}\) have equal derivatives through second order; a third
derivative remainder gives the one-step error
\(h^3B_k^3e^{hB_k}/3\).  Contractivity and a telescoping sum over \(n=t/h\)
steps prove Eq.~\eqref{eq:loop-kinetic-bound}.
\end{proof}

The response is operational.  Preparing
\(\rho_{\boldsymbol\sigma,\eta}
=2^{-N}(\mathbf1+\eta P_{\boldsymbol\sigma})\) and measuring
\(P_{\boldsymbol\tau}\) gives a matrix element of \(\Phi_h^n\).  In
particular,
\begin{align}
 R^Z_{j\leftarrow i}(t)
 &=\left[e^{-(\gamma/d)L_Gt}\right]_{ji},\nonumber\\
 R^X_{j\leftarrow i}(t)
 &=e^{-2\kappa t}
   \left[e^{-(\gamma/d)L_Gt}\right]_{ji},
 \label{eq:loop-local-response}
\end{align}
up to the controlled finite-step error
\eqref{eq:loop-kinetic-bound}.  On a ring, taking
\(t=N^2s/\gamma\) and a compatible \(h(N)\) yields the heat kernel on the
unit circle, with separately bounded product-formula and lattice errors.
This realizes the microscopic--kinetic--hydrodynamic separation advocated
in Sec.~\ref{sec:status-map}.

\paragraph{Boundary of the result.}
Fresh one-pass bath cells are a physical memory assumption.  Reusing the
same cell gives coherent recurrences, for example \(\cos(2n\theta)\), rather
than exponential decay.  The circuit also uses an external schedule and a
calibrated collision strength.  Equation~\eqref{eq:loop-kinetic-bound}
therefore derives the response kinetics under a controlled Markovian
dilation; it does not derive an autonomous clock, a thermal reservoir, or a
metric.

\subsection{Boolean charge records and metastable classical motion}
\label{sec:loop-boolean}

\paragraph{Assumption removed.}
The first finite open-system loop removes asymptotic abelianness as a
separate dynamical hypothesis.  For the Lindbladian
\(\mathcal L_G=\mathcal D+\mathcal A\) in
Eq.~\eqref{eq:loop-kinetic-generator}, the invariant algebra, its complete
gap, and the weakly perturbed sector dynamics are all calculable
\cite{XuBoolean2026}.  General fixed-algebra theory identifies stationary
operator structures of quantum dynamical semigroups
\cite{Frigerio1978}; the result below additionally fixes the complete
many-body gap and the topology dependence of the metastable window.

Let \(\mathcal E\) be the trace-preserving conditional expectation
\begin{equation}
 \mathcal E(X)=\sum_{m=0}^{N}
 \frac{\operatorname{Tr}(P_mX)}{\binom Nm}P_m .
 \label{eq:loop-boolean-expectation}
\end{equation}

\begin{theorem}[Noise-selected Boolean algebra and complete gap]
\label{thm:loop-boolean}
For connected \(d\)-regular \(G\),
\begin{align}
 \operatorname{Fix}(e^{t\mathcal L_G})
 &=\bigoplus_{m=0}^{N}\mathbb CP_m,\nonumber\\
 \lambda_G&=\min\!\left\{2\kappa,
 \frac{\gamma}{d}\lambda_2(L_G)\right\},
 \label{eq:loop-boolean-gap}\\
 \|[e^{t\mathcal L_G}(A),e^{t\mathcal L_G}(B)]\|_2
 &\le 2e^{-\lambda_Gt}
 \bigl[\|B\|_\infty\|A-\mathcal EA\|_2
 +\|A\|_\infty\|B-\mathcal EB\|_2\bigr].
 \label{eq:loop-boolean-comm}
\end{align}
The norm \(\|\cdot\|_2\) is the normalized Hilbert--Schmidt norm.
\end{theorem}

Local dephasing removes all occupation-basis off-diagonal matrix units.
On the diagonal subspace, exchange is the symmetric exclusion process.
Aldous' spectral-gap theorem identifies its gap with the one-particle graph
gap; a uniform one-excitation coherence attains the off-diagonal value
\(2\kappa\).  Thus neither the algebra nor the slow scale is inferred from a
small-system diagonalization.

The theorem preserves the full charge distribution:
\begin{equation}
 \rho_\infty=\sum_{m=0}^{N}
 p_m\,\frac{P_m}{\binom Nm},
 \qquad p_m=\operatorname{Tr}(P_m\rho).
 \label{eq:loop-boolean-state}
\end{equation}
All intrasection quantum information is erased, while the atoms \(P_m\)
remain exactly distinguishable.  The finite-time Boolean defect obeys
\begin{equation}
 \mathfrak b(t)=
 \sup_{\|A\|_\infty,\|B\|_\infty\le1}
 \|[e^{t\mathcal L_G}(A),e^{t\mathcal L_G}(B)]\|_2
 \le4e^{-\lambda_Gt}.
 \label{eq:loop-boolean-defect}
\end{equation}

Now perturb by
\(\mathcal L_{G,\epsilon}=\mathcal L_G+\epsilon\mathcal V\), where
\(\mathcal V\) generates a bistochastic Hilbert--Schmidt contraction, and
put \(v=\|\mathcal V\|_{2\to2}\),
\(\delta=\epsilon v/\lambda_G\).  The compressed operator
\(\mathcal K=\mathcal E\mathcal V\mathcal E\) is a backward classical Markov
generator on the atoms \(P_m\).  Standard spectral perturbation and a
Duhamel expansion give, on the metastable time window,
\begin{equation}
 \left\|
 e^{t\mathcal L_{G,\epsilon}}
 -e^{t\epsilon\mathcal K}\mathcal E
 \right\|_{2\to2}
 \le C\!\left(e^{-\lambda_Gt}+\delta
 +\epsilon^2 v^2t/\lambda_G\right).
 \label{eq:loop-metastable}
\end{equation}
Uniform bit flips preserve the sector manifold and give an exact Ehrenfest
chain.  A one-site flip does not commute with \(\mathcal E\), produces
genuine leakage, and is controlled by the same ratio \(\delta\).

Topology enters twice.  On rings
\(\lambda_2(L_G)=O(N^{-2})\), so a fixed isolation quality requires
\(\epsilon v/\gamma=O(N^{-2})\).  On bounded-degree expanders,
\(\lambda_2(L_G)=O(1)\), and a size-independent perturbation is tolerated.
The result therefore ties a graph property to the formation and robustness
of a classical information manifold.

\paragraph{Boundary of the result.}
The uniform theorem is in normalized Hilbert--Schmidt norm.  It controls
the stated operator-space defect and mixed-state families of bounded
normalized purity, but it is not a dimension-free diamond-norm or
worst-case local operator-norm theorem.  Nor does the model decide which
charge should encode a spacetime record.

\subsection{A transition between full quantum memory and Boolean records}
\label{sec:loop-memory}

\paragraph{Assumption removed.}
A fixed Boolean algebra does not show whether non-Abelian information and
classical records can be separated by a genuine dynamical phase transition.
This is realized by a contact-process flag field that gates local memory
dephasing \cite{XuMemory2026}.  The flag quotient is the classical contact
process \cite{Harris1974,Liggett1999}, while the object that crosses the
transition is an operational quantum memory channel.

Let \(n_i(t)\in\{0,1\}\) be contact-process occupation variables and let a
memory qudit at \(i\) dephase at rate \(\kappa\) only when \(n_i=1\).
After waiting to \(t_w\), an arbitrary state is written into a finite window
\(W\).  For matrix units
\(E_{\boldsymbol a\boldsymbol b}\), define
\(D(\boldsymbol a,\boldsymbol b)=\{i:a_i\ne b_i\}\).  The reduced channel
obeys the exact operator identity
\begin{align}
 \Phi^*_{s,t}(E_{\boldsymbol a\boldsymbol b})
 &=q_D(s,t)E_{\boldsymbol a\boldsymbol b},\nonumber\\
 q_D(s,t)&=\mathbb E\exp\!\left[
 -\kappa\int_s^t\!du\sum_{i\in D}n_i(u)\right].
 \label{eq:loop-memory-fk}
\end{align}
This is a Feynman--Kac representation of the entire memory channel, not a
formula for one chosen coherence.

Define the late-write algebra by taking the infinite-volume or
quasistationary flag limit before the unbounded storage interval:
\begin{equation}
 \mathcal A_{\rm tail}(W)=
 \left\{A:\lim_{t_w\to\infty}\sup_{\tau\ge0}
 \|\Phi^*_{t_w,t_w+\tau}(A)-A\|=0\right\}.
 \label{eq:loop-tail-definition}
\end{equation}

\begin{theorem}[Tail-algebra and capacity transition]
\label{thm:loop-memory}
For the contact process on \(\mathbb Z^r\), away from criticality,
\begin{equation}
 \mathcal A_{\rm tail}(W)=
 \begin{cases}
 M_{d^{|W|}},&\lambda<\lambda_c,\\
 \operatorname{diag}(M_{d^{|W|}})
 \simeq\mathbb C^{d^{|W|}},&\lambda>\lambda_c.
 \end{cases}
 \label{eq:loop-memory-algebra}
\end{equation}
Below threshold there are \(C,\alpha>0\) such that
\begin{equation}
 1-q_D(t_w,\infty)
 \le\frac{\kappa|D|C}{\alpha}e^{-\alpha t_w}.
 \label{eq:loop-memory-bound}
\end{equation}
For one memory qubit entangled with an untouched reference,
\begin{equation}
 \mathcal N_{R:M}=\frac q2,
 \qquad
 Q_{\rm mem}(q)=
 1-h_2\!\left(\frac{1-q}{2}\right)
 \label{eq:loop-memory-capacity}
\end{equation}
for independently reinitialized channel uses.
\end{theorem}

Subcritical sharpness makes the late integrated activity finite and proves
\eqref{eq:loop-memory-bound}.  In the upper invariant supercritical process,
ergodicity makes the occupation integral diverge for every nonempty \(D\),
so only diagonal matrix units survive.  The one-qubit map is a phase-flip
channel with \(p=(1-q)/2\); its Choi state has negativity \(q/2\), and
degradability gives Eq.~\eqref{eq:loop-memory-capacity}.  The same scalar
therefore controls noncommutativity, entanglement transmission, and
memoryless-use quantum capacity.

Quantum absorbing-state models can support coherence in the activity sector
itself and can depart from directed-percolation criticality
\cite{WamplerCooper2025}.  That is not the mechanism used here.  The flags
form a classical quotient; the additional result is the exact matrix-unit
Feynman--Kac action, which determines the full tail algebra and both
reference-entanglement and capacity witnesses without a moment closure.

\paragraph{Boundary of the result.}
On every finite connected graph the flags eventually absorb and all memories
become dark.  The phase transition belongs to the stated thermodynamic or
quasistationary order of limits.  The activity quotient is classical and
does not receive backaction from the memories.  A nonzero background
dephasing rate rounds the ideal infinite-time distinction into a lifetime
crossover.  These qualifications do not alter the exact channel identity,
but they delimit the phase claim.

\subsection{Causal order as an absorbing dissipative phase}
\label{sec:loop-causal}

\paragraph{Assumption removed.}
The earlier causal-order criterion assumed a vertex function \(T\) satisfying
\(u\to v\Rightarrow T(v)>T(u)\).  A local defect model produces the DAG
first and constructs \(T\) afterward \cite{XuCausalOrder2026}.

The underlying undirected graph is a chain, or more generally a locally
finite tree, of triangular blocks.  Each edge carries a qubit whose monitored
record is an orientation.  In block \(r\), the two cyclic words \(000\) and
\(111\) define
\begin{equation}
 C_r=|000\rangle\langle000|_r+|111\rangle\langle111|_r,
 \qquad Q_r=\mathbf1-C_r .
 \label{eq:loop-cycle-projector}
\end{equation}
Local Lindblad jumps heal a cyclic block at total rate \(g\), while a cyclic
neighbor infects a transitive target at rate \(1/2\).  The jump set is
covariant under global arrow reversal, so no global direction is encoded in
the generator.

\begin{theorem}[Order from defect extinction]
\label{thm:loop-causal-order}
For every recorded orientation \(\omega\) of a block cactus, the threshold
response graph is acyclic if and only if \(C_r(\omega)=0\) for every block.
When this holds,
\begin{equation}
 T_\omega(v)=\max_{p:\,p\to v}|p|
 \label{eq:loop-longest-path}
\end{equation}
is finite, obeys \(u\to v\Rightarrow T_\omega(v)>T_\omega(u)\), and the
reachability relation is a locally finite strict partial order.
\end{theorem}

The block-cut tree is essential: every simple undirected cycle belongs to one
block.  Hence the local projectors certify all directed cycles.  On a generic
lattice a triangle test can miss a longer chordless loop.

The commuting algebra generated by \(\{C_r\}\) is invariant and has the exact
classical restriction
\begin{align}
 (\mathcal K_gf)(\boldsymbol n)
 ={}&g\sum_r n_r
 [f(\boldsymbol n^{r,0})-f(\boldsymbol n)]\nonumber\\
 &+\frac12\sum_r(1-n_r)(n_{r-1}+n_{r+1})
 [f(\boldsymbol n^{r,1})-f(\boldsymbol n)] .
 \label{eq:loop-contact-generator}
\end{align}
After rescaling time, this is the one-dimensional contact process with
\(\lambda=1/g\) \cite{Harris1974,Liggett1999}.  Thus
\begin{equation}
 g_c=\lambda_c^{-1}=0.3032280(18)
 \label{eq:loop-causal-gc}
\end{equation}
using the accepted numerical contact-process threshold.  In the ordered
phase, every fixed operational window \(W\) satisfies
\begin{equation}
 1-P_{\rm DAG}(W,t)
 \le A(g)|W|e^{-\alpha(g)t},
 \qquad
 P_{\rm DAG}(W,t)=
 \operatorname{Tr}\!\left[\rho(t)\prod_{r\in W}Q_r\right].
 \label{eq:loop-causal-window}
\end{equation}
The vanishing rate \(\alpha(g)\) is an order-formation gap, and the exact
quotient transfers the directed-percolation critical exponents to cycle
density and causal correlation scales.

\paragraph{Boundary of the result.}
The model assumes the block-cactus adjacency, the edge registers, the local
laboratory runtime, monitoring, and a thresholded response readout.  It
generates neither the undirected graph nor a volume measure.  Its theorem
removes an assumed Lyapunov ranking on a nontrivial graph family; it is not
yet a background-free spacetime theorem.

\subsection{Quantum-field response as an inverse Lorentzian observable}
\label{sec:loop-lorentz}

\paragraph{Assumption removed.}
The first inverse-geometry model encoded proper time into a calibrated
depolarizing parameter.  That construction has been replaced by a physical
field response: a Ramsey ratio isolates a Wightman covariance, whose
rank defect determines an unknown de Sitter radius
\cite{XuLorentz2026}.  Scalar propagators and localized quantum probes are
known to carry metric and curvature information
\cite{SaravaniAslanbeigiKempf2016,PercheMartinMartinez2022}, while
Lorentzian multidimensional scaling embeds supplied interval data
\cite{CloughEvans2017}.  The distinct claim tested here is partial-data
completion from a physical covariance kernel: anchor responses fix an
unknown curvature scale and predict target pairs never used in the fit.

For real spacelike field smearings \(\Phi_i,\Phi_j\), two qubit probes undergo
\begin{equation}
 U=\exp[-ig(\sigma_z^i\Phi_i+\sigma_z^j\Phi_j)] .
 \label{eq:loop-ramsey-unitary}
\end{equation}
In a zero-mean quasifree state their same- and opposite-sign coherences are
\(\chi_\pm=\exp\{-2g^2[V_i+V_j\pm2W_{ij}]\}\), so
\begin{equation}
 W_{ij}=\frac{1}{8g^2}\log\frac{\chi_-}{\chi_+}.
 \label{eq:loop-ramsey-ratio}
\end{equation}
The local self-variances cancel exactly.

For a conformally coupled massless scalar in the Bunch--Davies state on
\(\mathrm{dS}_4(R)\),
\begin{equation}
 W_{ij}=\frac{1}{8\pi^2R^2(1-Z_{ij})}
 \label{eq:loop-ds-kernel}
\end{equation}
on the spacelike branch.  For six generic anchors define
\(A_{ii}=0\), \(A_{ij}=(8\pi^2W_{ij})^{-1}\).

\begin{theorem}[Finite-anchor curvature and blind completion]
\label{thm:loop-lorentz}
If \(A\) is invertible and the six ambient anchor vectors span
\(\mathbb R^{1,4}\), then
\begin{equation}
 R^2=(\mathbf1^TA^{-1}\mathbf1)^{-1}.
 \label{eq:loop-radius}
\end{equation}
The candidate is admissible only if
\(G=R^2\mathbf1\mathbf1^T-A\) has rank five and Lorentzian inertia
\((-++++0)\).  With \(K\ge6\) anchors and
\(b_x=R^2\mathbf1-A_{Ax}\), every withheld target--target invariant is
\begin{equation}
 \widehat Z_{xy}=\frac{b_x^TG^+b_y}{R^2}.
 \label{eq:loop-blind-completion}
\end{equation}
\end{theorem}

The determinant lemma applied to
\(G=X\eta X^T\) proves Eq.~\eqref{eq:loop-radius}; the pseudoinverse identity
proves Eq.~\eqref{eq:loop-blind-completion}.  Redundant anchor sextuples,
Lorentzian inertia, target hyperboloid norms, and unopened target pairs
overidentify the model.  Perturbation bounds propagate Ramsey shot noise,
finite smearing, and Gram conditioning into errors on \(R\) and
\(\widehat Z_{xy}\).  The blind pairs can therefore falsify the assumed
geometry--state kernel rather than merely reproduce data used in a fit.

\paragraph{Boundary of the result.}
The field species, Bunch--Davies state, de Sitter model class, spacelike
support, canonical normalization, probe coupling, and a nondegenerate anchor
frame are inputs.  Equation~\eqref{eq:loop-ramsey-ratio} is physical, but the
kernel \eqref{eq:loop-ds-kernel} supplies the geometric model.  The theorem is
an inverse problem inside that class; it does not derive a Lorentzian
manifold from the causal-order model or determine temporal orientation of a
timelike held-out pair.

\subsection{A fixed physical modular band}
\label{sec:loop-modular}

\paragraph{Assumption removed.}
Earlier lattice tests retained a fixed number of modular modes.  As the
cutoff is removed, that prescription narrows the physical band and can make
convergence tautological.  A state-independent band selected by the local
lattice-boost operator removes this loophole
\cite{XuModular2026}.

For the half-filled free-fermion interval, let \(C_N\) be the restricted
correlation matrix and
\begin{equation}
 h_N=\log[(\mathbf1-C_N)C_N^{-1}]
 \label{eq:loop-modular-h}
\end{equation}
the exact one-particle modular Hamiltonian.  The known tridiagonal commuting
operator determines a local boost matrix \(B_N\), with
\([B_N,h_N]=0\).  This construction is rooted in the
Bisognano--Wichmann form and its lattice implementations
\cite{BisognanoWichmann1975,Giudici2018}.  The operational band
\begin{equation}
 Q_\Lambda=\mathbf1_{[-\Lambda,\Lambda]}(B_N)
 \label{eq:loop-modular-band}
\end{equation}
is fixed before \(h_N\) is diagonalized; its rank grows with \(N\) at fixed
\(\Lambda\).

For any real Borel function \(g\), put \(G_N=g(B_N)\) and
\(E_N=h_N-G_N\).  Every Hermitian correlation perturbation satisfies
\begin{align}
 |\operatorname{Tr}(\delta C E_N)|
 \le{}&\delta_\Lambda[G_N]\,
 \|Q_\Lambda\delta C Q_\Lambda\|_1\nonumber\\
 &+\delta_\Lambda^\perp[G_N]\,
 \|\bar Q_\Lambda\delta C\bar Q_\Lambda\|_1,
 \label{eq:loop-modular-certificate}\\
 \delta_\Lambda[G_N]
 &=\|Q_\Lambda E_NQ_\Lambda\|_\infty,\qquad
 \delta_\Lambda^\perp[G_N]
 =\|\bar Q_\Lambda E_N\bar Q_\Lambda\|_\infty .
\end{align}
The in-band constant is minimax optimal for trace-neutral Gaussian tangents.
It therefore certifies all linear modular responses in the stated physical
band, rather than only one prepared wave packet.

Common-spectrum analysis gives the sparse correction
\begin{equation}
 \widetilde B_N=B_N+\frac1{N^2}
 \left[-\frac14B_N+\frac{1}{4\pi^2}B_N^3\right].
 \label{eq:loop-range-three}
\end{equation}
It contains only first- and third-neighbor hopping.  The cubic coefficient is
analytic; the linear coefficient was selected on \(N\le128\), frozen, and
tested without refitting on \(N=192,\ldots,768\).  For \(\Lambda=8\), the
archived deterministic data give
\begin{equation}
 \delta_8[B_N]\propto N^{-1.992},
 \qquad
 \delta_8[\widetilde B_N]\propto N^{-3.984},
 \label{eq:loop-modular-scaling}
\end{equation}
with mode-entry-aware fourth-order envelopes on the held-out sizes.  The
first relation is an analytic finite-band response certificate; the
\(N^{-4}\) corrected law is controlled numerical evidence, not an all-size
theorem.

\paragraph{Boundary of the result.}
This loop concerns a one-dimensional free state and a modular-to-lattice-
boost response.  It contains no transverse area, multiple null directions,
Newton constant, focusing equation, conservation completion, or dynamical
metric.  Consequently it is not an Einstein derivation.  Interacting and
higher-dimensional versions require new operator estimates.

\subsection{Dynamical charge center and selective fixed-sector response}
\label{sec:loop-center}

\paragraph{Assumption removed.}
A central variable should not be attached by hand.  A local
number-conserving channel derives a complete additive-charge fixed
center, its formation gap, and its lifting under weak symmetry breaking
\cite{XuCentral2026}.

On a connected graph let \(n_i\) be hard-core occupations,
\(\widehat V=\sum_i n_i\), and use local dephasing together with symmetric
number-conserving hops \(L_{ij}=\sigma_i^-\sigma_j^+\).  The reversible
Lindbladian has Dirichlet form
\begin{align}
 -\langle X,\mathcal L^*(X)\rangle
 ={}&\frac{\gamma}{2}\sum_i\|[n_i,X]\|_2^2\nonumber\\
 &+\frac{\eta}{2}\sum_{\{i,j\}\in E}
 \left(\|[L_{ij},X]\|_2^2+\|[L_{ji},X]\|_2^2\right).
 \label{eq:loop-center-dirichlet}
\end{align}

\begin{theorem}[Dynamical additive-charge center]
\label{thm:loop-center}
Let \(P_v\) project onto \(\widehat V=v\).  Then
\begin{align}
 \ker\mathcal L^*&=\operatorname{span}\{P_0,\ldots,P_N\},\nonumber\\
 \mathcal E_V(X)&=\sum_{v=0}^{N}
 \frac{\operatorname{Tr}(P_vX)}{\binom Nv}P_v,\nonumber\\
 \|e^{t\mathcal L^*}X-\mathcal E_VX\|_2
 &\le e^{-\Delta t}\|X-\mathcal E_VX\|_2,\qquad
 \Delta\ge\min\{\gamma/2,\eta\lambda_2(L_G)\}.
 \label{eq:loop-center-gap}
\end{align}
The one-point density profile closes exactly:
\begin{equation}
 \mathcal L^*(n_i)=\eta\sum_{j\sim i}(n_j-n_i).
 \label{eq:loop-center-heat}
\end{equation}
\end{theorem}

The fixed projections are dynamically produced, not an appended register.
They are central in the fixed algebra and for charge-preserving
laboratories, but not in the microscopic algebra
\(\mathcal B(\mathcal H)\).  Homogeneous weak flips give an exact Ehrenfest
chain on \(v=0,\ldots,N\); in the balanced stationary state,
\begin{equation}
 \langle\delta\widehat V(t)\delta\widehat V(0)\rangle
 =\frac N4e^{-2\epsilon t}.
 \label{eq:loop-center-correlation}
\end{equation}
This separates the graph-controlled formation time from the
symmetry-breaking charge-memory time.

There is a precise, but conditional, field-theory consequence.  If \(v\) is
identified with fixed Euclidean metric volume and the sector response is
normalized as \(Z_v[J]/Z_v[0]\), a pure volume counterterm \(c v\) cancels:
\begin{equation}
 \frac{Z_v[J;c]}{Z_v[0;c]}
 =\frac{e^{-cv}Z_v[J;0]}{e^{-cv}Z_v[0;0]}
 =\frac{Z_v[J;0]}{Z_v[0;0]} .
 \label{eq:loop-fixed-volume-cancellation}
\end{equation}
A coupling to an operator \(Q\) that is not constant inside the sector does
not cancel; its derivative is the connected covariance
\(-\operatorname{Cov}(O,Q)\).  Thus the statement is selective: it removes
only the identity or pure-volume part from conditional response and leaves
curvature-dependent counterterms to ordinary effective-field-theory
renormalization.  Fixed-total-volume and unimodular formulations provide
the closest gravitational precedent \cite{FiolGarriga2010}, while global
sequestering uses a different set of constraints \cite{KaloperPadilla}.
The microscopic result here is the local charge-center theorem; identifying
that charge with metric volume remains a separate bridge.

\paragraph{Boundary of the result.}
The microscopic \(U(1)\) law is assumed.  No theorem identifies
\(\widehat V\) with metric four-volume, derives sectorwise gravitational
normalization, predicts a cosmological constant, or explains why the
appropriate laboratory is confined to one volume sector.  Equation
\eqref{eq:loop-fixed-volume-cancellation} is therefore a conditional
interface attached to an exact charge-center theorem, not a solution of the
cosmological-constant problem.

\subsection{What these loops do and do not compose}
\label{sec:loop-noncomposition}

The completed loops remove several assumptions independently:
\begin{equation}
\begin{array}{rcl}
\text{postulated UCP kinetics}
&\Longrightarrow&\text{fresh-cell unitary dilation},\\
\text{assumed asymptotic abelianness}
&\Longrightarrow&\text{exact Boolean algebra and gap},\\
\text{static algebraic dichotomy}
&\Longrightarrow&\text{memory-capacity transition},\\
\text{assumed Lyapunov ranking}
&\Longrightarrow&\text{DAG followed by derived rank},\\
\text{calibrated proper-time oracle}
&\Longrightarrow&\text{field-covariance inverse geometry},\\
\text{shrinking fixed-mode audit}
&\Longrightarrow&\text{fixed physical modular band},\\
\text{attached central label}
&\Longrightarrow&\text{dynamical additive-charge center}.
\end{array}
\label{eq:loop-removals}
\end{equation}
No arrow in Eq.~\eqref{eq:loop-removals} licenses the replacement
\[
\text{collision graph}\equiv\text{causal graph}\equiv
\text{de Sitter anchors}\equiv\text{fermion interval}\equiv
\text{metric-volume sectors}.
\]
Those identifications would be additional assumptions.  A future unified
model must define common microscopic degrees of freedom and prove that its
distinct scaling regimes reduce to the models above, with an error budget
that remains controlled across both arrows
\[
\text{microscopic process}\longrightarrow\text{kinetic response}
\longrightarrow\text{order--volume geometry}.
\]
This noncomposition statement is as important as the positive theorems: it
prevents evidence obtained in one solvable model from being counted twice as
support for a different bridge.
